% ============================================================
%  IEEE-Style Project Report — SO-100 Robot Arm RL
%  Compile with:  pdflatex report.tex  (run twice for references)
%  Or:            latexmk -pdf report.tex
% ============================================================

\documentclass[conference]{IEEEtran}

% ── Packages ────────────────────────────────────────────────
\usepackage[T1]{fontenc}
\usepackage[utf8]{inputenc}
\usepackage{cite}                  % IEEE-style numbered citations [1]
\usepackage{amsmath,amssymb}       % math symbols and equations
\usepackage{graphicx}              % \includegraphics
\usepackage{booktabs}              % professional tables (\toprule etc.)
\usepackage{array}                 % extended column definitions in tables
\usepackage{tabularx}              % tables that fill text width
\usepackage{xcolor}                % colored text (optional)
\usepackage{url}                   % \url{} for links
\usepackage{hyperref}              % clickable links in PDF
\usepackage{caption}               % fine-grained caption control
\usepackage{subcaption}            % side-by-side figures with sub-captions
\usepackage{microtype}             % better text justification

% ── Page numbering ──────────────────────────────────────────
% IEEEtran suppresses page numbers by default; restore them here
\pagestyle{plain}                  % page number centred at bottom

% ── Hyperref setup ──────────────────────────────────────────
\hypersetup{
    colorlinks = true,
    linkcolor  = black,
    citecolor  = black,
    urlcolor   = blue,
    pdftitle   = {Reinforcement Learning for SO-100 Robot Arm Manipulation},
    pdfauthor  = {Your Name},
}


% ── Custom commands (add your own as needed) ─────────────────
\newcommand{\todo}[1]{\textcolor{red}{[TODO: #1]}}  % red TODO markers


% ============================================================
%  TITLE & AUTHORS
% ============================================================
\begin{document}

\title{Reinforcement Learning for Robotic Manipulation:\\
       Training the SO-100 Arm to Touch and Place Objects}

\author{
    \IEEEauthorblockN{Your Full Name}
    \IEEEauthorblockA{
        Department of \todo{Your Department}\\
        \todo{Your University}\\
        \todo{your.email@university.edu}
    }
}

\maketitle

% ── Running footer note (optional) ──────────────────────────
\begin{abstract}
%% ── ABSTRACT ──────────────────────────────────────────────
%%  150–250 words. State: what you did, how, key results.
%%  Do NOT cite references in the abstract.

This paper presents the design, implementation, and evaluation of a
reinforcement learning (RL) system for controlling the SO-100, a six
degree-of-freedom robotic arm, in a simulated MuJoCo environment.
The objective is to train an autonomous agent—using the Soft
Actor-Critic (SAC) algorithm—to perform manipulation tasks of
increasing difficulty, starting with touching a target cube and
progressing toward pick-and-place operations.
The environment is built on top of the \texttt{gym-so100} package,
which wraps dm\_control physics into a Gymnasium-compatible interface.
Three task variants are defined: dense-reward touch, sparse-reward
touch, and cube-to-bin placement.
A 12-dimensional state observation is used, comprising object
positions and joint angles.
Training is conducted for 300\,000 environment steps using
parallel vectorized environments and observation normalisation.
Results show that the agent successfully learns an approach
strategy—reducing end-effector-to-cube distance from over 0.18~m to
under 0.03~m—but does not yet achieve reliable contact closure
within the tested training budget.
Qualitative analysis of the reward function reveals that the
step penalty and contact-bonus structure present conflicting
incentives, which are discussed alongside directions for future work.

\end{abstract}

% IEEE keywords
\begin{IEEEkeywords}
reinforcement learning, robotic manipulation, Soft Actor-Critic,
MuJoCo, reward shaping, SO-100 robot arm, Gymnasium
\end{IEEEkeywords}


% ============================================================
\section{Introduction}
\label{sec:introduction}
% ============================================================

Robotic manipulation is one of the most challenging open problems in
autonomous systems research.
Even tasks that are trivial for humans—reaching for an object,
grasping it, and moving it to a target location—require a robot to
solve a high-dimensional continuous control problem under physical
constraints and sensory uncertainty \cite{todorov2012mujoco}.

Reinforcement learning offers a promising framework for robot
control: rather than hand-coding a controller for every possible
configuration, the agent learns a policy by interacting with the
environment and maximising cumulative reward \cite{sutton2018rl}.
Recent advances in off-policy algorithms, particularly
Soft Actor-Critic (SAC) \cite{haarnoja2018sac}, have demonstrated
impressive results on continuous-action robotic benchmarks.

This project applies these ideas to the SO-100 robot arm
\cite{gymso100}, a six-joint manipulator originally designed for
teleoperation tasks.
Using MuJoCo as the physics backend and the \texttt{gym-so100}
package as the environment interface, we train an SAC agent to
perform three manipulation tasks of increasing complexity.

\subsection{Motivation}

The SO-100 arm is a low-cost, open-source manipulator platform.
Developing RL-based controllers for it bridges the gap between
academic RL research and practical, affordable robotics.
This project is intended to serve as a reproducible baseline
demonstrating what 300\,000 training steps of SAC can achieve on
this hardware.

\subsection{Contributions}

The main contributions of this work are:
\begin{itemize}
    \item A fully configured RL training pipeline for the SO-100
          arm using SAC and Stable Baselines 3.
    \item A reward function analysis identifying the cause of
          incomplete contact learning.
    \item An interactive visualisation and reward-probe toolkit
          for diagnosing agent behaviour.
    \item Recommendations for reward redesign and extended training.
\end{itemize}


% ============================================================
\section{Problem Statement}
\label{sec:problem}
% ============================================================

\subsection{Task Definition}

We define the manipulation problem as a Markov Decision Process
(MDP) $\mathcal{M} = (\mathcal{S}, \mathcal{A}, \mathcal{R},
\mathcal{P}, \gamma)$, where:

\begin{itemize}
    \item $\mathcal{S} \subset \mathbb{R}^{12}$ is the state space,
    \item $\mathcal{A} = [-1,1]^6$ is the continuous action space,
    \item $\mathcal{R}: \mathcal{S} \times \mathcal{A} \to \mathbb{R}$
          is the reward function defined in Section~\ref{sec:reward},
    \item $\mathcal{P}$ is the transition dynamics governed by
          MuJoCo physics at 50~Hz,
    \item $\gamma = 0.99$ is the discount factor.
\end{itemize}

\subsection{State Observation}

At each timestep $t$, the agent receives a 12-dimensional
observation vector:
\begin{equation}
    \mathbf{s}_t = \bigl[
        \mathbf{p}_\text{cube},\;
        \mathbf{p}_\text{bin},\;
        \mathbf{p}_\text{ee},\;
        \mathbf{q}
    \bigr] \in \mathbb{R}^{12}
    \label{eq:obs}
\end{equation}
where $\mathbf{p}_\text{cube}, \mathbf{p}_\text{bin},
\mathbf{p}_\text{ee} \in \mathbb{R}^3$ are the 3-D Cartesian
positions of the cube, bin centre, and end-effector respectively,
and $\mathbf{q} \in \mathbb{R}^6$ contains the current joint angles.

\subsection{Action Space}

The action $\mathbf{a}_t \in [-1,1]^6$ is a normalised joint
position command.
Before being sent to the simulator, each component is
un-normalised to the real joint range:

\begin{table}[h]
    \centering
    \caption{Joint Ranges and Normalisation}
    \label{tab:joints}
    \begin{tabular}{llrr}
        \toprule
        Index & Joint & Min (rad) & Max (rad) \\
        \midrule
        0 & Waist          & $-1.92$ & $+1.92$ \\
        1 & Shoulder       & $-3.32$ & $+0.174$ \\
        2 & Elbow          & $-0.174$ & $+3.14$ \\
        3 & Forearm Roll   & $-1.66$ & $+1.66$ \\
        4 & Wrist Rotate   & $-2.79$ & $+2.79$ \\
        5 & Gripper        & $-0.174$ & $+1.75$ \\
        \bottomrule
    \end{tabular}
\end{table}

\subsection{Episode Structure}

Each episode begins with the arm at a fixed home pose and the cube
spawned at a uniformly random position within:
$x \in [-0.25, -0.15]$~m, $y \in [0.30, 0.60]$~m, $z = 0.05$~m.
An episode terminates either on success (reward = 4) or after
300 steps ($T = 6$~seconds of simulated time).


% ============================================================
\section{Approaches and Solutions}
\label{sec:approach}
% ============================================================

\subsection{Environment Architecture}

The environment stack is illustrated in Fig.~\ref{fig:env_stack}.

%% ── HOW TO INSERT A FIGURE ──────────────────────────────────
%%  1. Place your image in report/figures/env_stack.png
%%  2. Uncomment the block below
%%
% \begin{figure}[h]
%     \centering
%     \includegraphics[width=\columnwidth]{figures/env_stack.png}
%     \caption{Environment architecture: MuJoCo physics wrapped by
%              dm\_control, then by SO100Env (Gymnasium), then by
%              VecNormalize and SubprocVecEnv for parallel training.}
%     \label{fig:env_stack}
% \end{figure}

The layered design separates concerns cleanly:
\begin{enumerate}
    \item \textbf{MuJoCo / dm\_control} simulates rigid-body physics
          at 50~Hz.
    \item \textbf{SO100Env} (Gymnasium wrapper) exposes the
          standard \texttt{reset()} / \texttt{step()} interface and
          formats observations.
    \item \textbf{SubprocVecEnv} runs $N=4$ independent environment
          copies in separate processes, multiplying experience throughput.
    \item \textbf{VecNormalize} maintains a running mean and standard
          deviation of observations and normalises them to
          approximately $\mathcal{N}(0,1)$ before the policy network
          sees them.
\end{enumerate}

\subsection{Reward Function}
\label{sec:reward}

The reward function used in the \texttt{SO100TouchCubeTask} is a
shaped dense reward consisting of four components:

\subsubsection{Distance Shaping}

Let $d_t = \|\mathbf{p}_\text{ee} - \mathbf{p}_\text{cube}\|_2$
be the end-effector-to-cube Euclidean distance at step $t$.
The shaping term is:
\begin{equation}
    r_\text{shape}(d) =
    \max \left(
        0.1\!\left(1-\tfrac{d}{0.7}\right),\;
        0.2\!\left(1-\tfrac{d}{0.5}\right),\;
        0.5\!\left(1-\tfrac{d}{0.3}\right),\;
        1.0\!\left(1-\tfrac{d}{0.1}\right),\;
        2.0\!\left(1-\tfrac{d}{0.05}\right)
    \right)
    \label{eq:shape}
\end{equation}
saturated at $r_\text{shape} = 0$ when $d \geq 0.7$.

\subsubsection{Contact Bonus}

A binary contact bonus is awarded when any gripper pad geometry
makes physical contact with the cube geometry:
\begin{equation}
    r_\text{contact} = \mathbf{1}[\text{touch\_gripper}] \times 1.0
\end{equation}

\subsubsection{Step Penalty}

A constant penalty of $-0.2$ is applied every timestep to
encourage faster task completion.

\subsubsection{Success Return}

If the gripper is touching the cube and $d < 0.05$~m simultaneously,
the episode is terminated and a terminal reward of $+4.0$ is issued.

The total step reward is therefore:
\begin{equation}
    r_t =
    \begin{cases}
        4.0 & \text{if touch\_gripper} \wedge d < 0.05 \\
        r_\text{shape}(d_t) + r_\text{contact} - 0.2 & \text{otherwise}
    \end{cases}
    \label{eq:reward}
\end{equation}

\subsection{Policy: Soft Actor-Critic}

SAC \cite{haarnoja2018sac} is an off-policy actor-critic algorithm
that maximises a maximum-entropy objective:
\begin{equation}
    J(\pi) = \sum_{t=0}^{T}
    \mathbb{E}_{(\mathbf{s}_t,\mathbf{a}_t)\sim\rho_\pi}
    \!\left[ r(\mathbf{s}_t, \mathbf{a}_t)
    + \alpha\,\mathcal{H}(\pi(\cdot|\mathbf{s}_t)) \right]
\end{equation}
where $\alpha$ is the automatically-tuned entropy coefficient and
$\mathcal{H}$ is the policy entropy.

Both the actor and twin-critic networks are fully connected MLPs
with architecture $12 \to 256 \to 256 \to \cdot$
(Table~\ref{tab:hyperparams}).

\begin{table}[h]
    \centering
    \caption{SAC Hyperparameters}
    \label{tab:hyperparams}
    \begin{tabular}{ll}
        \toprule
        Hyperparameter & Value \\
        \midrule
        Network architecture   & $[256, 256]$ MLP \\
        Learning rate          & $3 \times 10^{-4}$ \\
        Replay buffer size     & 100\,000 \\
        Batch size             & 256 \\
        Entropy coeff.\ $\alpha$ & auto-tuned \\
        Discount $\gamma$      & 0.99 \\
        Learning starts        & 1\,000 steps \\
        Parallel environments  & 4 \\
        Total training steps   & 300\,000 \\
        Observation type       & State (12-dim) \\
        \bottomrule
    \end{tabular}
\end{table}


% ============================================================
\section{Testing and Validation}
\label{sec:testing}
% ============================================================

\subsection{Training Curves}

%% ── HOW TO INSERT SIDE-BY-SIDE FIGURES ──────────────────────
%%  Place images in report/figures/ and uncomment:
%%
% \begin{figure}[h]
%     \centering
%     \begin{subfigure}[b]{0.48\columnwidth}
%         \includegraphics[width=\textwidth]{figures/ep_rew_mean.png}
%         \caption{Mean episode reward}
%         \label{fig:rew}
%     \end{subfigure}
%     \hfill
%     \begin{subfigure}[b]{0.48\columnwidth}
%         \includegraphics[width=\textwidth]{figures/ent_coef.png}
%         \caption{Entropy coefficient}
%         \label{fig:ent}
%     \end{subfigure}
%     \caption{TensorBoard training curves over 300\,000 steps.}
%     \label{fig:curves}
% \end{figure}

Training was monitored via TensorBoard.
The mean episode reward increased steadily from approximately
$-60$ at the start of training (random policy) to a peak of
\todo{insert peak value} around step \todo{insert step}.
The entropy coefficient $\alpha$ decreased monotonically from
$\approx 1.0$ to $\approx 0.2$, indicating a transition from
broad exploration to more deterministic exploitation.

\subsection{Quantitative Evaluation}

After training, the final policy was evaluated over 10 episodes
with deterministic action selection.
Results are summarised in Table~\ref{tab:eval}.

\begin{table}[h]
    \centering
    \caption{Evaluation Results (300\,000 steps, deterministic policy)}
    \label{tab:eval}
    \begin{tabular}{lrr}
        \toprule
        Metric & Mean & Std \\
        \midrule
        Episode reward          & \todo{val} & \todo{val} \\
        Episode length (steps)  & \todo{val} & \todo{val} \\
        Success rate (\%)       & \todo{val} & —          \\
        Min EE-cube distance (m)& \todo{val} & \todo{val} \\
        \bottomrule
    \end{tabular}
\end{table}

\subsection{Qualitative Behaviour}

Trajectory analysis of the deterministic policy across three
episodes revealed the following behaviour pattern:

\begin{enumerate}
    \item The agent consistently reduces the end-effector-to-cube
          distance from $\approx 0.19$~m to $\approx 0.03$~m within
          the first 15--40 steps.
    \item The arm then oscillates around the cube without closing
          the gripper, collecting the distance-shaping reward
          indefinitely.
    \item No contact events (\texttt{touch\_gripper = True}) are
          observed in any of the three evaluation episodes.
\end{enumerate}

\subsection{Reward Probe Analysis}

An interactive reward-probe tool was developed to inspect the
reward signal in real time during manual control.
This confirmed that the distance-shaping component provides a
strong, smooth gradient toward the cube, while the contact bonus
remains latent — the agent never discovers it within the training
budget because random exploration rarely produces a gripper-closing
motion precisely when already within $0.05$~m of the cube.

\subsection{Checkpoint Comparison}

\todo{Compare 3 checkpoints (e.g.\ 30k, 150k, 300k steps) and
describe the improvement in behaviour between them.}


% ============================================================
\section{Future Aspects}
\label{sec:future}
% ============================================================

\subsection{Reward Redesign}

The current reward function contains two conflicting pressures
when $d < 0.05$~m: the step penalty ($-0.2$) discourages staying
near the cube, while the contact bonus ($+1.0$) only fires on
actual physical contact.
Proposed modifications include:
\begin{itemize}
    \item Adding an explicit gripper-closing reward term when
          $d < 0.05$~m, proportional to the gripper joint position.
    \item Increasing the contact bonus from $+1.0$ to $+3.0$ to
          make it a dominant signal worth exploring for.
    \item Removing or reducing the step penalty to allow the agent
          more time near the cube without negative reinforcement.
\end{itemize}

\subsection{Extended Training}

Training for 1\,000\,000 steps (approximately $3\times$ longer)
may allow the agent to stumble upon contact events during
exploration and subsequently reinforce this behaviour.
GPU-accelerated training would make this feasible within a
reasonable time budget.

\subsection{Curriculum Learning}

A curriculum could progressively increase task difficulty:
spawn the cube directly under the gripper initially, then
gradually increase randomness in the cube's initial position.
This would allow the agent to first learn contact closure, then
learn to approach from a distance.

\subsection{Harder Tasks}

Once contact is mastered, the natural progression is:
\begin{enumerate}
    \item \textbf{Lift:} reward the agent for raising the cube
          above the table.
    \item \textbf{Place:} reward the agent for depositing the
          cube in the target bin (\texttt{SO100CubeToBin-v0}).
    \item \textbf{Pixel observations:} replace the 12-number
          state vector with a camera image for a more realistic
          and challenging setting.
\end{enumerate}

\subsection{Sim-to-Real Transfer}

The ultimate goal is to transfer the learned policy from simulation
to the physical SO-100 arm.
Domain randomisation (varying object masses, friction, and cube
positions during training) and a more faithful sensor model would
improve transfer fidelity.


% ============================================================
\section{Conclusion}
\label{sec:conclusion}
% ============================================================

This project demonstrated the end-to-end application of the
Soft Actor-Critic algorithm to robotic manipulation of the SO-100
arm in a MuJoCo simulation environment.
A complete training and evaluation pipeline was built from scratch,
including environment configuration, reward function design,
parallel training, checkpoint management, and interactive
visualisation tools.

The trained agent exhibits learned approach behaviour—reliably
reducing the end-effector-to-cube distance to within 3~cm—but does
not yet achieve consistent contact, indicating that 300\,000 steps
with the current reward structure is insufficient to bridge the
final gap between approach and grasp.

Reward function analysis identified the root cause: the contact
bonus is too weak relative to the distance shaping, and the step
penalty discourages the persistent near-object exploration needed
to discover the contact signal.
Targeted reward modifications and extended training are expected
to resolve this in future work.


% ============================================================
%  REFERENCES
% ============================================================
\bibliographystyle{IEEEtran}

%% ── HOW TO USE REFERENCES ───────────────────────────────────
%%  Option A (BibTeX — recommended):
%%    1. Create report/references.bib with your entries
%%    2. Uncomment the line below:
% \bibliography{references}

%%  Option B (inline — no .bib file needed):
%%    Keep the thebibliography block below as-is and add entries.

\begin{thebibliography}{9}

\bibitem{todorov2012mujoco}
E.~Todorov, T.~Erez, and Y.~Tassa,
``MuJoCo: A physics engine for model-based control,''
in \emph{Proc.\ IEEE/RSJ Int.\ Conf.\ Intelligent Robots and Systems
(IROS)}, 2012, pp.~5026--5033.

\bibitem{sutton2018rl}
R.~S.~Sutton and A.~G.~Barto,
\emph{Reinforcement Learning: An Introduction}, 2nd~ed.
Cambridge, MA: MIT Press, 2018.

\bibitem{haarnoja2018sac}
T.~Haarnoja, A.~Zhou, P.~Abbeel, and S.~Levine,
``Soft Actor-Critic: Off-policy maximum entropy deep reinforcement
learning with a stochastic actor,''
in \emph{Proc.\ Int.\ Conf.\ Machine Learning (ICML)}, 2018,
pp.~1861--1870.

\bibitem{gymso100}
I.~Julczuk,
``gym-so100-c: A Gymnasium environment for the SO-100 robot arm,''
2024. [Online]. Available: \url{https://github.com/ilonajulczuk/gym-so100-c}

\bibitem{sb3}
A.~Raffin \emph{et al.},
``Stable-Baselines3: Reliable Reinforcement Learning Implementations,''
\emph{J.\ Machine Learning Research}, vol.~22, no.~268, pp.~1--8, 2021.

\bibitem{dmcontrol}
S.~Tunyasuvunakool \emph{et al.},
``dm\_control: Software package for physics-based simulation and
reinforcement learning,''
\emph{Software Impacts}, vol.~6, p.~100022, 2020.

%% Add more references here following the same \bibitem pattern.
%% Example:
% \bibitem{yourkey}
% A.~Author and B.~Author,
% ``Title of the paper,''
% \emph{Journal Name}, vol.~X, no.~Y, pp.~Z1--Z2, Year.

\end{thebibliography}

\end{document}
